% ============================================================
%  03_marco_teorico.tex — Capítulo 3: Marco teórico
% ============================================================

\chapter{Marco teórico}
\label{cap:marco-teorico}

\textit{[Aquí se explican los conceptos, teorías y tecnologías que sustentan la propuesta.]}

\section{Invernaderos y variables ambientales}
\label{sec:invernaderos}

\textit{[Tipos de invernadero, variables (temperatura, humedad, CO\textsubscript{2}, luz,
riego), y su importancia para el cultivo.]}

\section{Instrumentación y sensórica}
\label{sec:sensorica}

\textit{[Sensores y actuadores, acondicionamiento de señales y adquisición de datos.]}

\section{Internet de las cosas (IoT)}
\label{sec:iot}

\textit{[Arquitecturas IoT, protocolos (MQTT, CoAP, HTTP), puertas de enlace y la nube.]}

\section{Ontologías y Web Semántica}
\label{sec:web-semantica}

\begin{description}
    \item[Ontología:] definición, componentes (clases, propiedades, individuos, axiomas).
    \item[Lenguajes:] RDF, RDFS, OWL 2.
    \item[Herramientas:] Protégé, razonadores (HermiT, Pellet), SPARQL.
\end{description}

\section{Modelos de aprendizaje automático}
\label{sec:ml}

\textit{[Algoritmos de regresión, clasificación, series temporales y su aplicación al
control climático.]}

\section{Metodología de construcción de ontologías}
\label{sec:metodologias-ontologia}

\textit{[Metodologías como METHONTOLOGY, NeOn, u otras; evaluación de ontologías.]}
